\section{Input sensitivity and interpretation}
\label{sec:sensitivity}
\subsection{Source loading and thermodynamic information}
The Q2 cases expose a sensitivity hidden by prescribing the iodide
split. \Cref{fig:thermo} shows the BiI$_3$ iodine fraction
\begin{equation}
 f_{\rm BiI_3}=
 \frac{3N^{\rm wall}_{\rm BiI_3}}
 {2N^{\rm wall}_{\rm PbI_2}+3N^{\rm wall}_{\rm BiI_3}}.
 \label{eq:split}
\end{equation}
This includes only iodine in the two modelled deposited iodides.
It is distinct from the BiI$_3$ fraction of all experimental iodine,
which also contains KI. Additional metal-saturation cases fill selected
points between the original 0.1\% and 1\% branches; line segments
connect computed points and are not a fitted release law.

At 1\% metal saturation on the fine steel mesh, the baseline
thermodynamics gives $f_{\rm BiI_3}=94.57\%$. Raising the BiI(g)
standard Gibbs energy by 29 kJ/mol gives 2.00\% with the same numerical
resolution. The large difference survives refinement, demonstrating
conditional thermodynamic sensitivity rather than numerical noise.
The perturbation is applied consistently to formation and reactive-wall
constants. A constant Gibbs-energy shift scales the formation constant
by $\exp[-29{,}000/(RT)]$, approximately 0.051 at 1173 K and
0.00093 at 500 K. It is a substantial surrogate perturbation, not an
independently measured correction to BiI thermodynamics.

Source loading also matters. With baseline thermodynamics the steel
BiI$_3$ share at 0.2\%, 0.316\% and 0.562\% metal saturation is
52.50\%, 70.58\% and 94.56\%; the shifted branch gives 47.66\%,
37.43\% and 15.88\%. The low-loading trend is not universally
monotonic. The original steel 0.1\% baseline branch gives 56.47\%
and an 81.10\% scan overlap, making it a useful diagnostic comparison,
but not an identified physical saturation. The published deposits alone
cannot uniquely distinguish source loading, thermodynamic information
and omitted chemistry. Choosing a branch because it best matches one
scan would introduce calibration and require independent validation.

\begin{figure}[tb]
\centering
\includegraphics[width=.96\textwidth]{source_thermodynamics.pdf}
\caption{Conditional source-loading and BiI thermodynamic sensitivity
for steel and silica LBE. Symbols denote completed Q2 cases; segments
only connect those points. Dashed grey lines show the published
BiI$_3$ share after excluding KI \cite{Liu2025chemical}.
The $+29$ kJ/mol branch is a surrogate perturbation, not a second
qualified thermodynamic dataset.}
\label{fig:thermo}
\end{figure}

\subsection{Silica temperature-coordinate discrepancy}
\label{sec:silica}
The published silica scan places PbI$_2$ and BiI$_3$ maxima at 51 and
59 cm, labelled 306$\pm$24 and 127$\pm$39\,$^\circ$C. At those
coordinates, the digitised temperature curve gives approximately 369
and 235\,$^\circ$C: differences of 63 and 108 K. On that same
curve, 306 and 127\,$^\circ$C occur at 55.096 and 63.653 cm.
The discrepancy exceeds the nominal raster temperature resolution and
suggests a temperature-coordinate consistency problem in the supplied
inputs. It does not identify which measurement, annotation or registration
requires correction.

Two diagnostic mappings shift only the silica temperature curve upstream
by 4.096 or 4.653 cm, aligning one of those curve crossings with the
corresponding labelled position. Equivalently,
$T_\delta(x)=T_0(x+\delta)$ for $\delta>0$, with
$\Delta x_T=-\delta$. The gamma scan remains fixed. These shifts are
derived from the published curve and labels, not optimised against the
CFD profiles. All original unshifted results are retained.

\Cref{tab:alignment,fig:alignment} shows that the shifts move Q1 maxima
from 56/63 cm to 52/59 or 51/58 cm and improve normalised overlap from
7.20\% to 75.29\% or 81.00\%. The low-loading Q2 variants place
both maxima at 51/59 cm, with overlaps 79.16\% and 83.33\%.
This explains why a plausible coordinate change has a larger effect
than the tested source-position or time-step changes. It does not
confirm either shift as the experimental truth. The two shifts also
produce different peak temperatures and cannot simultaneously reconcile
both original temperature labels by a uniquely established mapping.
\begin{table}[tb]
\centering
\caption{Silica temperature-coordinate hypotheses. Shifts and peak
positions are in cm; raw L1 is in percentage points and overlap in
percent. Q2 uses baseline thermodynamics and 0.1\% metal saturation.
Every row uses the same fixed experimental scan.}
\label{tab:alignment}
\small\begin{tabular}{lrrrrr}
\toprule
Source & $\Delta x_T$ & $x_{\rm PbI_2}$ & $x_{\rm BiI_3}$ & Raw L1 & Overlap \\
\midrule
Q1 & 0.00 & 56 & 63 & 181.84 & 7.20 \\
Q1 & -4.10 & 52 & 59 & 46.28 & 75.29 \\
Q1 & -4.65 & 51 & 58 & 37.05 & 81.00 \\
Q2 & 0.00 & 56 & 63 & 181.84 & 7.20 \\
Q2 & -4.10 & 51 & 59 & 39.40 & 79.16 \\
Q2 & -4.65 & 51 & 59 & 34.31 & 83.33 \\
\bottomrule
\end{tabular}

\end{table}
\begin{figure}[tb]
\centering
\includegraphics[width=\textwidth]{silica_alignment.pdf}
\caption{Input-derived silica mappings and their consequences.
Top: the original reconstructed temperature curve, two shifts and
printed peak labels. Bottom: common-scan-normalised iodine
profiles for Q1 and low-loading Q2. Improved agreement is conditional
on an unconfirmed coordinate hypothesis. Published data are adapted
from Liu et al.\ \cite{Liu2025chemical}.}
\label{fig:alignment}
\end{figure}

The shifted low-loading Q2 cases retain a BiI$_3$ share near 23.97\%,
about five percentage points below the published iodide-only split.
Correct peak positions therefore do not ensure the correct partition.
Resolving the discrepancy requires the original thermocouple profile
and its coordinate registration, rather than selecting the CFD variant
with the highest overlap.

\subsection{Transport, kinetics, flow and shutdown}
\Cref{tab:sensitivity} reports the range of iodine profile changes from
the designated baseline within each tested category, using
\cref{eq:numeric-l1}. They are responses to finite hypothetical changes,
not statistical confidence limits. The source-loading row uses its
1\% reference branch; the temperature-coordinate row measures shifts
from the original silica input and can approach 200\% when profile
supports move apart. These baseline-relative differences must not be
confused with raw discrepancies against the experimental scan.
\begin{table}[tb]
\centering
\caption{Ranges of baseline-relative iodine inventory-profile changes
(percent) over each completed sensitivity category. Shutdown compares
80 s inventories with their 60 s restart state; other rows compare
60 s inventories. Case counts and variations are in \cref{tab:study}.}
\label{tab:sensitivity}
\small\begin{tabular}{lrr}
\toprule
Assumption & Minimum & Maximum \\
\midrule
Gas diffusion & 0.15 & 14.20 \\
Iodide accommodation & 0.25 & 4.21 \\
Source position & 0.08 & 0.19 \\
Q2 metal loading & 0.00 & 112.14 \\
Flow reference & 5.09 & 10.93 \\
Silica temperature coordinates & 184.68 & 192.50 \\
Flow-stop relaxation & 1.10 & 4.63 \\
\bottomrule
\end{tabular}

\end{table}

Diffusion alternatives change profiles by up to 14.20\%, and a
273.15 K flow reference raises inferred helium mass flow by 9.15\%
relative to 298.15 K and changes profiles by up to 10.93\%.
Iodide accommodation changes give up to 4.21\%, whereas moving the
prescribed source region by $\pm1.5$ cm gives at most 0.19\%.
This small source-position response concerns the selected empty-tube
source model; it cannot rule out an effect of a real boat on flow or
thermal fields. The maxima for these categories are shown in
\cref{fig:physical} and should be compared with the unresolved
2.06--10.58\% fine/medium spatial differences. Where numerical and
input changes have comparable magnitude, the present results do not
rank their physical importance precisely.

The four fixed-temperature flow-stop cases change deposited iodine
profiles by 1.10--4.63\% over 20 s, with centroid shifts no larger
than 0.0422 cm. They show that residual gas can alter wall inventories
after release stops. The column is neither cooling nor supplied with
a clean-helium flush in these runs. Consequently, they do not predict
the actual end-of-experiment redistribution after the furnaces are
turned off. Likewise, equilibrium speciation and the irreversible
BiI wall channel remain unqualified kinetic approximations.

\begin{figure}[tb]
\centering
\includegraphics[width=.8\textwidth]{physical_sensitivity.pdf}
\caption{Maximum baseline-relative iodine profile changes among the
selected physical-input comparisons. Flow-stop relaxation uses an
80 s endpoint; the others use 60 s. These maxima describe the tested
case set, not a general uncertainty budget or an experimentally
identified parameter ranking.}
\label{fig:physical}
\end{figure}
